\par\noindent\begin{minipage}{\linewidth}
% !TeX root = ../教材题库_全部条目_转写工作稿.tex
% Source: XJ-B1-C01-PR-P006-Y0346
% 第1章 / 1.1 集合 / 教材 p.6 / PDF 第13页
% 已对照原页转写题干及公式；待教师复核；未制作答案
% 原题及相邻原解截图（核对用）：../题目原图/XJ-B1-C01-PR-P006-Y0346-01.webp
% 本文件只在转写完成后作为题干使用；OCR 条目尚不是可靠题干。
\begin{bankitem}{p.6 练习 第2题}{第1章 / 1.1 集合 / 教材 p.6 / PDF 第13页}{XJ-B1-C01-PR-P006-Y0346}{已对照原页转写题干及公式；待教师复核；未制作答案}
用描述法表示下列集合：
\subq{1}$\{1,3,5,7,9\}$；
\subq{2}不等式 $3x>2$ 的解集；
\subq{3}$\{3,5,7,11,13,17,19\}$。
\end{bankitem}
\end{minipage}\par\medskip
